% =============================================================================
% ВСПОМОГАТЕЛЬНЫЕ ЭЛЕМЕНТЫ
% =============================================================================

\tcbuselibrary{theorems, skins, breakable, xparse, minted}

\newcommand{\blockMainText}[1]{%
    \ignorespaces#1\par\normalfont
}

\newcommand{\customtitlefont}{
    \sffamily\bfseries\large
}

\newcommand{\boxsepline}{
    \par\vspace{3mm plus 1mm minus 1mm}
    {\color{black!40}\hrule height 1.4pt}
    \vspace{2.5mm plus 1mm minus 1mm}
    \noindent
}

\newcommand{\printRemark}[1]{%
    \if\relax\detokenize{#1}\relax\else
        \boxsepline
        {\small
            \settowidth{\noteiconwidth}{\noteicon}%
            \setlength{\parindent}{\noteiconwidth}%
            \noindent\noteicon\itshape\color{black!80}#1\par
        }%
    \fi
}

\newcommand{\printProof}[1]{%
    \if\relax\detokenize{#1}\relax\else
        \boxsepline
        \begingroup
            \pushQED{\qed}%
            \textbf{Доказательство.}\ #1\unskip\popQED\par
        \endgroup
    \fi%
}

\newcommand{\printExamples}[1]{%
    \if\relax\detokenize{#1}\relax\else
        \boxsepline
        \begin{itemize}[leftmargin=*, label=\textcolor{listL1}{\textbullet}, itemsep=0.5em, parsep=0pt, topsep=3pt]
            #1
        \end{itemize}
    \fi
}

\newcommand{\printSolution}[1]{%
  \if\relax\detokenize{#1}\relax\else
    \boxsepline
    \begingroup
        \pushQED{\qed}%
        \textbf{Решение.}\ #1\unskip\popQED\par
    \endgroup
  \fi
}

\newcommand{\printOutput}[1]{%
    \if\relax\detokenize{#1}\relax\else
        \boxsepline
        \textbf{Вывод.}\ #1\par
    \fi
}

\newcommand{\printComplexity}[1]{%
    \if\relax\detokenize{#1}\relax\else
        \boxsepline
        \textbf{Сложность.}\ #1\par
    \fi
}


% =============================================================================
% НАСТРОЙКА СЧЕТЧИКОВ И СТИЛЯ
% =============================================================================

\newcounter{commonThmCounter}[section]
\renewcommand{\thecommonThmCounter}{\thesection.\arabic{commonThmCounter}}

\tcbset{
  myboxstyle/.style={
    enhanced jigsaw,
    breakable,
    enforce breakable,
    colback=white,
    colframe=#1,
    coltext=black,
    colbacktitle=#1,
    coltitle=white,
    fonttitle=\customtitlefont,
    parbox=false,
    arc=3mm,
    boxrule=0.7mm,
    opacityback=1,
    top=2.5mm, bottom=2mm, left=2mm, right=2mm,
    before skip=6mm plus 2mm minus 1mm,
    after skip=6mm plus 2mm minus 1mm,
    pad at break*=3mm,
    before upper app={%
      \setlength{\parindent}{1.5em}%
      \setlength{\parskip}{0.4em plus 0.25em minus 0.2em}%
      \clubpenalty=2000%
      \widowpenalty=2000%
      \displaywidowpenalty=2000%
      \brokenpenalty=0%
      \interlinepenalty=0%
      \noindent\ignorespaces
    },
  },
}

\numberwithin{equation}{chapter}


% =============================================================================
% ВСПОМОГАТЕЛЬНЫЕ МАКРОСЫ ДЛЯ ЗАГОЛОВКОВ И МЕТОК
% =============================================================================

\newcommand{\blockSetLabel}[1]{\ifstrempty{#1}{}{\label{#1}}}

\newlength{\blockTitleHang}
\newsavebox{\blockPrefixBox}

\newif\ifblockBadgeNames
\blockBadgeNamestrue

\newcommand{\blockBadge}[1]{%
    \tikz[baseline=(B.base)]
        \node[
            draw=white,
            line width=1pt,
            rounded corners=2.6pt,
            fill=white,
            text=tcbcolframe,
            inner xsep=6pt,
            inner ysep=2.5pt,
            minimum width=2em,
            anchor=base,
            font=\sffamily\bfseries\Large,
        ] (B) {\strut #1};%
}

\newcommand{\blockName}[2]{%
    \ifblockBadgeNames\blockBadge{#1}\else #2\fi
}

\newcommand{\blockTitle}[2]{%
    \sbox{\blockPrefixBox}{#1\nobreakspace\hspace{0.25em}\thecommonThmCounter\ifstrempty{#2}{}{.}}%
    \settowidth{\blockTitleHang}{\usebox{\blockPrefixBox}\space}%
    \hangindent=\blockTitleHang
    \hangafter=1
    \hyphenpenalty=10000\exhyphenpenalty=10000
    \pretolerance=9999\tolerance=2000
    \hbadness=10000
    \raggedright
    \usebox{\blockPrefixBox}\ifstrempty{#2}{}{ #2}%
}

\newcommand{\blockTitleStar}[2]{%
    \sbox{\blockPrefixBox}{#1}%
    \settowidth{\blockTitleHang}{\usebox{\blockPrefixBox}\hspace{0.5em}}%
    \hangindent=\blockTitleHang
    \hangafter=1
    \hyphenpenalty=10000\exhyphenpenalty=10000
    \pretolerance=9999\tolerance=2000
    \hbadness=10000
    \raggedright
    \usebox{\blockPrefixBox}\ifstrempty{#2}{}{\hspace{0.5em}#2}%
}


% =============================================================================
% ОПРЕДЕЛЕНИЕ ОКРУЖЕНИЙ (НУМЕРОВАННЫХ)
% =============================================================================

\newtcolorbox{defiBox}[2]{
  myboxstyle=deficolor,
  code={\refstepcounter{commonThmCounter}\blockSetLabel{#1}},
  title={\blockTitle{\blockName{О}{Определение}}{#2}},
}

\newtcolorbox{statBox}[2]{
    myboxstyle=statcolor,
    code={\refstepcounter{commonThmCounter}\blockSetLabel{#1}},
    title={\blockTitle{\blockName{У}{Утверждение}}{#2}},
}

\newtcolorbox{lemmBox}[2]{
    myboxstyle=lemmcolor,
    code={\refstepcounter{commonThmCounter}\blockSetLabel{#1}},
    title={\blockTitle{\blockName{Л}{Лемма}}{#2}},
}

\newtcolorbox{theoBox}[2]{
    myboxstyle=theocolor,
    code={\refstepcounter{commonThmCounter}\blockSetLabel{#1}},
    title={\blockTitle{\blockName{Т}{Теорема}}{#2}},
}

\newtcolorbox{concBox}[2]{
    myboxstyle=conccolor,
    code={\refstepcounter{commonThmCounter}\blockSetLabel{#1}},
    title={\blockTitle{\blockName{С}{Следствие}}{#2}},
}

\newtcolorbox{examBox}[2]{
    myboxstyle=examcolor,
    code={\refstepcounter{commonThmCounter}\blockSetLabel{#1}},
    title={\blockTitle{\blockName{П}{Пример}}{#2}},
}

\newtcolorbox{noteBox}[2]{
    myboxstyle=notecolor,
    code={\refstepcounter{commonThmCounter}\blockSetLabel{#1}},
    title={\blockTitle{\blockName{З}{Замечание}}{#2}},
}

\newtcolorbox{algoBox}[2]{
    myboxstyle=algocolor,
    code={\refstepcounter{commonThmCounter}\blockSetLabel{#1}},
    title={\blockTitle{\blockName{А}{Алгоритм}}{#2}},
}

\newtcolorbox[auto counter, number within=chapter]{taskBox}{
    enhanced,
    colback=white,
    colframe=white!50!black,
    coltitle=black,
    fonttitle=\bfseries,
    detach title,
    before upper={\textbf{\tcbtitle.}\ },
    frame hidden,
    left=0mm, right=0mm, top=1mm, bottom=1mm,
    title={\thetcbcounter}
}


% =============================================================================
% ОПРЕДЕЛЕНИЕ ОКРУЖЕНИЙ (НЕНУМЕРОВАННЫХ)
% =============================================================================

\newtcolorbox{defiBoxStar}[2]{
    myboxstyle=deficolor,
    title={\blockTitleStar{\blockName{О}{Определение}}{#2}},
}

\newtcolorbox{statBoxStar}[2]{
    myboxstyle=statcolor,
    title={\blockTitleStar{\blockName{У}{Утверждение}}{#2}},
}

\newtcolorbox{lemmBoxStar}[2]{
    myboxstyle=lemmcolor,
    title={\blockTitleStar{\blockName{Л}{Лемма}}{#2}},
}

\newtcolorbox{theoBoxStar}[2]{
    myboxstyle=theocolor,
    title={\blockTitleStar{\blockName{Т}{Теорема}}{#2}},
}

\newtcolorbox{concBoxStar}[2]{
    myboxstyle=conccolor,
    title={\blockTitleStar{\blockName{С}{Следствие}}{#2}},
}

\newtcolorbox{examBoxStar}[2]{
    myboxstyle=examcolor,
    title={\blockTitleStar{\blockName{П}{Пример}}{#2}},
}

\newtcolorbox{noteBoxStar}[2]{
    myboxstyle=notecolor,
    title={\blockTitleStar{\blockName{З}{Замечание}}{#2}},
}

\newtcolorbox{algoBoxStar}[2]{
    myboxstyle=algocolor,
    title={\blockTitleStar{\blockName{А}{Алгоритм}}{#2}},
}


% =============================================================================
% КОМАНДЫ-ОБЕРТКИ
% =============================================================================

% \defi{Метка}{Термин}{Текст}[Примеры][Замечание]
\NewDocumentCommand{\defi}{ s m m +m +O{} +O{} }{
    \IfBooleanTF{#1}{%
        \begin{defiBoxStar}{#2}{#3}
            \blockMainText{#4}
            \par \normalfont
            \printExamples{#5}
            \printRemark{#6}
        \end{defiBoxStar}
    }{%
        \begin{defiBox}{#2}{#3}
            \blockMainText{#4}
            \par \normalfont
            \printExamples{#5}
            \printRemark{#6}
        \end{defiBox}
    }
}

% \stat{Метка}{Название}{Текст}[Док-во][Замечание]
\NewDocumentCommand{\stat}{ s m m +m +O{} +O{} }{
    \IfBooleanTF{#1}{%
        \begin{statBoxStar}{#2}{#3}
            \blockMainText{#4}
            \par \normalfont
            \printProof{#5}
            \printRemark{#6}
        \end{statBoxStar}
    }{%
        \begin{statBox}{#2}{#3}
            \blockMainText{#4}
            \par \normalfont
            \printProof{#5}
            \printRemark{#6}
        \end{statBox}
    }
}

% \lemm{Метка}{Название}{Текст}[Док-во][Замечание]
\NewDocumentCommand{\lemm}{ s m m +m +O{} +O{} }{
    \IfBooleanTF{#1}{%
        \begin{lemmBoxStar}{#2}{#3}
            \blockMainText{#4}
            \par \normalfont
            \printProof{#5}
            \printRemark{#6}
        \end{lemmBoxStar}
    }{%
        \begin{lemmBox}{#2}{#3}
            \blockMainText{#4}
            \par \normalfont
            \printProof{#5}
            \printRemark{#6}
        \end{lemmBox}
    }
}

% \theo{Метка}{Название}{Текст}[Док-во][Замечание]
\NewDocumentCommand{\theo}{ s m m +m +O{} +O{} }{
    \IfBooleanTF{#1}{%
        \begin{theoBoxStar}{#2}{#3}
            \blockMainText{#4}
            \par \normalfont
            \printProof{#5}
            \printRemark{#6}
        \end{theoBoxStar}
    }{%
        \begin{theoBox}{#2}{#3}
            \blockMainText{#4}
            \par \normalfont
            \printProof{#5}
            \printRemark{#6}
        \end{theoBox}
    }
}

% \conc{Метка}{Название}{Текст}[Док-во][Замечание]
\NewDocumentCommand{\conc}{ s m m +m +O{} +O{} }{
    \IfBooleanTF{#1}{%
        \begin{concBoxStar}{#2}{#3}
            \blockMainText{#4}
            \par \normalfont
            \printProof{#5}
            \printRemark{#6}
        \end{concBoxStar}
    }{%
        \begin{concBox}{#2}{#3}
            \blockMainText{#4}
            \par \normalfont
            \printProof{#5}
            \printRemark{#6}
        \end{concBox}
    }
}

% \exam{Метка}{Название}{Текст / Условие}[Решение][Замечание]
\NewDocumentCommand{\exam}{ s m m +m +O{} +O{} }{
    \IfBooleanTF{#1}{%
        \begin{examBoxStar}{#2}{#3}
            \blockMainText{#4}
            \par \normalfont
            \printSolution{#5}
            \printRemark{#6}
        \end{examBoxStar}
    }{%
        \begin{examBox}{#2}{#3}
            \blockMainText{#4}
            \par \normalfont
            \printSolution{#5}
            \printRemark{#6}
        \end{examBox}
    }
}

% \note{Метка}{Название}{Текст}
\NewDocumentCommand{\note}{ s m m +m }{
    \IfBooleanTF{#1}{%
        \begin{noteBoxStar}{#2}{#3}
            \blockMainText{#4}
        \end{noteBoxStar}
    }{%
        \begin{noteBox}{#2}{#3}
            \blockMainText{#4}
        \end{noteBox}
    }
}

% \algo{Метка}{Название}{Описание алгоритма}[Сложность][Замечание]
\NewDocumentCommand{\algo}{ s m m +m +O{} +O{} }{
    \IfBooleanTF{#1}{%
        \begin{algoBoxStar}{#2}{#3}
            \blockMainText{#4}
            \par \normalfont
            \printComplexity{#5}
            \printRemark{#6}
        \end{algoBoxStar}
    }{%
        \begin{algoBox}{#2}{#3}
            \blockMainText{#4}
            \par \normalfont
            \printComplexity{#5}
            \printRemark{#6}
        \end{algoBox}
    }
}

% \task{Текст}
\NewDocumentCommand{\task}{ +m }{
    \begin{taskBox}
        \blockMainText{#1}
    \end{taskBox}
}


% =============================================================================
% КОД И ВЫВОД ПРОГРАММ
% =============================================================================

% --- Загрузка Lua-препроцессора и автоматическая чистка ----------------------- % ←
\directlua{code_filter = dofile("style/code_filter.lua")}                          % ←
\AtEndDocument{\directlua{code_filter.cleanup()}}                                  % ←

% Иконка-маркер стандартного потока вывода
\newcommand{\stdoutIcon}{%
    \raisebox{0.05ex}{\faAngleRight}%
    \kern0.1em
    \rule[-0.1ex]{0.45em}{0.55pt}%
}

% Бейдж: иконка + номер в круглых скобках                                         % ←
\newcommand{\stdoutBadge}[1]{%
    \tikz[baseline=(B.base)]
        \node[
            rounded corners=2pt,
            draw=xmatrixDark!55,
            line width=0.35pt,
            fill=xmatrix!8,
            text=xmatrixDark,
            inner xsep=5pt,
            inner ysep=1.2pt,
            font=\sffamily\bfseries\scriptsize,
        ] (B) {\strut\stdoutIcon\hspace{0.45em}#1};%
}

\newcommand{\cmo}[1]{%
    \unskip
    \nobreak\null\hfill\nobreak
    \mbox{\stdoutBadge{#1}}%
}

% Параметры левой колонки вывода программы
\newlength{\codeOutputNumberWidth}
\newlength{\codeOutputNumberGap}
\newlength{\codeOutputRuleWidth}
\newlength{\codeOutputTextGap}
\newlength{\codeOutputLeftMargin}
\newlength{\codeOutputNumberSep}

\setlength{\codeOutputNumberWidth}{1.0em} % колонка номеров (под иконкой)
\setlength{\codeOutputNumberGap}{0.55em}  % зазор номер -> линия
\setlength{\codeOutputRuleWidth}{0.4pt}   % толщина линии
\setlength{\codeOutputTextGap}{0.75em}    % зазор линия -> текст

\setlength{\codeOutputLeftMargin}{%
    \dimexpr
        \codeOutputNumberWidth
        + \codeOutputNumberGap
        + \codeOutputRuleWidth
        + \codeOutputTextGap
    \relax
}

% numbersep = расстояние от ПРАВОГО края номера до левого края текста
\setlength{\codeOutputNumberSep}{%
    \dimexpr
        \codeOutputNumberGap
        + \codeOutputRuleWidth
        + \codeOutputTextGap
    \relax
}

% Стиль номеров строк блока вывода программы
\newcommand{\codeOutputLineNumberStyle}{%
    \renewcommand{\theFancyVerbLine}{%
        \makebox[\codeOutputNumberWidth][c]{%
            \ttfamily\scriptsize\color{xmatrixDark!70}%
            \arabic{FancyVerbLine}%
        }%
    }%
}

% Шапка: иконка слева, затем текст                                                % ←
\tcbset{
  codeOutputBoxStyle/.style={
    enhanced,
    breakable,
    colback=xmatrix!4,
    colframe=xmatrixDark!55,
    colbacktitle=xmatrixDark!85,
    coltitle=white,
    fonttitle=\sffamily\bfseries\scriptsize,
    title={%
        \makebox[0pt][l]{\stdoutIcon}%
        \hspace*{\codeOutputLeftMargin}%
        \textsc{Вывод программы}\hfill
    },
    boxrule=0.3pt,
    arc=1.4mm,
    left=1mm, right=1mm, top=1mm, bottom=1mm,
    toptitle=0.5mm, bottomtitle=0.5mm,
    titlerule=0pt,
    before skip=2.2mm plus 0.6mm minus 0.4mm,
    after skip=0pt,
    parbox=false,
    % Непрерывная вертикальная линия-разделитель
    % (interior.north west / south west — границы внутренней области без рамки)
    overlay={%
    \draw[xmatrixDark!25, line width=\codeOutputRuleWidth]
        ([xshift=\dimexpr\codeOutputNumberWidth+\codeOutputNumberGap+0.5\codeOutputRuleWidth\relax]
            interior.north west)
        --
        ([xshift=\dimexpr\codeOutputNumberWidth+\codeOutputNumberGap+0.5\codeOutputRuleWidth\relax]
            interior.south west);
    },
  },
}

% Печать содержимого файла вывода: номера слева, в скобках                        % ←
\newcommand{\printCodeOutputFile}[1]{%
    \if\relax\detokenize{#1}\relax\else
        \par\nobreak
        \begingroup
            \codeOutputLineNumberStyle
            \IfFileExists{#1}{%
                \begin{tcolorbox}[codeOutputBoxStyle]
                    \inputminted[
                        fontsize=\small,
                        breaklines=true,
                        breakanywhere=true,
                        autogobble=false,
                        tabsize=4,
                        linenos=true,
                        numbers=left,
                        numbersep=\codeOutputNumberSep,   % ← было \codeOutputTextGap
                        xleftmargin=\codeOutputLeftMargin,
                        xrightmargin=0.4em,
                        numberblanklines=true,
                    ]{text}{#1}%
                \end{tcolorbox}%
            }{%
                \begin{tcolorbox}[codeOutputBoxStyle]
                    \noindent
                    \textcolor{xred}{%
                        \sffamily\bfseries\small
                        Файл вывода не найден: \texttt{\detokenize{#1}}%
                    }%
                \end{tcolorbox}%
            }%
        \endgroup
    \fi
}

% \code{Метка}{Название}{Файл с кодом}[Файл с выводом]
\NewTCBInputListing{\code}{ m m m O{} }{
    myboxstyle=codecolor,
    bottom=1mm,
    after skip=4mm plus 1mm minus 1mm,
    code={%
        \refstepcounter{commonThmCounter}%
        \blockSetLabel{#1}%
        % препроцессинг исходного файла → создаём временный *.stripped            % ←
        \directlua{code_filter.process(                                            % ←
            "\luaescapestring{#3}",                                                % ←
            "\luaescapestring{#3.stripped}")}%                                     % ←
    },
    title={\blockTitle{\blockName{К}{Код}}{#2}},
    listing engine=minted,
    listing only,
    listing file={#3.stripped},                                                    % ← .stripped
    minted language=python,
    minted options={
        fontsize=\small,
        breaklines=true,
        breakanywhere=false,
        autogobble=true,
        tabsize=4,
        escapeinside=«»,
    },
    after upper app={\printCodeOutputFile{#4}},
}